%%%%%%%%%%%%%%%%%%%%%%% file template.tex %%%%%%%%%%%%%%%%%%%%%%%%%
%
% This is a general template file for the LaTeX package SVJour3
% for Springer journals.          Springer Heidelberg 2010/09/16
%
% Copy it to a new file with a new name and use it as the basis
% for your article. Delete % signs as needed.
%
% This template includes a few options for different layouts and
% content for various journals. Please consult a previous issue of
% your journal as needed.
%
%%%%%%%%%%%%%%%%%%%%%%%%%%%%%%%%%%%%%%%%%%%%%%%%%%%%%%%%%%%%%%%%%%%
%
% First comes an example EPS file -- just ignore it and
% proceed on the \documentclass line
% your LaTeX will extract the file if required
%\begin{filecontents*}{example.eps}
%!PS-Adobe-3.0 EPSF-3.0
%%BoundingBox: 19 19 221 221
%%CreationDate: Mon Sep 29 1997
%%Creator: programmed by hand (JK)
%%EndComments
%gsave
%newpath
%  20 20 moveto
%  20 220 lineto
%  220 220 lineto
%  220 20 lineto
%closepath
%2 setlinewidth
%gsave
%  .4 setgray fill
%grestore
%stroke
%grestore
%\end{filecontents*}
%\RequirePackage{fix-cm}
%
\documentclass{svjour3}                     % onecolumn (standard format)
%\documentclass[smallcondensed]{svjour3}     % onecolumn (ditto)
%\documentclass[smallextended]{svjour3}       % onecolumn (second format)
%\documentclass[twocolumn]{svjour3}          % twocolumn
%
\smartqed  % flush right qed marks, e.g. at end of proof
%

%\setlength{\parskip}{1.4em}
\usepackage{mathptmx}
\usepackage{amssymb}
\usepackage{subfig}
%\usepackage{mathtools}
\usepackage{amsmath}
\usepackage{pgf,tikz}
\usepackage{float}
\usepackage{placeins}   % \FloatBarrier: floats stay in their subsection
% float placement: let LaTeX fill the pages (no half-empty pages before large figures and tables)
\renewcommand{\topfraction}{0.9}\renewcommand{\bottomfraction}{0.8}\renewcommand{\textfraction}{0.07}
\renewcommand{\floatpagefraction}{0.75}\setcounter{topnumber}{3}\setcounter{bottomnumber}{2}\setcounter{totalnumber}{5}
\extrafloats{100}
% subcaption is not loaded: it is incompatible with subfig (\subfloat), which is used for the sub-figures
\usepackage{algorithmic}
%\usepackage{floatpag}
\usepackage{url}
\usepackage{rotating}
\usepackage{algorithm}
\usepackage{multicol}
\usepackage{multirow}
%\usepackage{pdfpages}
\usepackage[justification=centering]{caption}
\usepackage{mathrsfs}
\usetikzlibrary{arrows}
\usepackage{array}
\usepackage{color}
\usepackage{caption}
%\usepackage{afterpage}
%\usepackage{tkz-tab}
\usepackage[justification=centering]{caption}
\usepackage{tikz}
\usepackage{adjustbox}
%\usepackage{cite}
%\usepackage{ulem}
%\usepackage{breakurl}
%\usepackage{bigints}
\usepackage{booktabs}
\usepackage{tabularx}
\usepackage{caption}
\usepackage{xcolor}
\newtheorem{theo}{Theorem}
\newtheorem{lem}{Lemma}
\newtheorem{cor}{Corrolary}
\newtheorem{defi}{Definition}
\usepackage{lineno}
\usepackage{longtable}
\newtheorem{examp}{Example}
\newtheorem{prop}{Proposition}
%\usepackage{indentfirst}
\usepackage{array}
\newcolumntype{P}[1]{>{\centering\arraybackslash}p{#1}}
\newcolumntype{M}[1]{>{\centering\arraybackslash}m{#1}}


\usepackage{enumerate}
\newtheorem{rem}{Remarks}
%\newtheorem{claim}{Claim}
\newtheorem{procedure}{Procedure}

\linespread{1.3}
\hoffset=0pt
\voffset=-3pt
\headsep=30pt
\setlength{\oddsidemargin}{0in}
\setlength{\evensidemargin}{-0.3in}
\setlength{\textwidth}{16cm}
%\setlength{\topmargin}{-0.3in}
\setlength{\textheight}{21truecm }
%\setcounter{MaxMatrixCols}{20}
%\marginparwidth=105pt
%\marginparsep=11pt

%
% \usepackage{mathptmx}      % use Times fonts if available on your TeX system
%
% insert here the call for the packages your document requires
%\usepackage{latexsym}
% etc.
%
% please place your own definitions here and don't use \def but
% \newcommand{}{}
%
% Insert the name of "your journal" with
% \journalname{myjournal}
%

% ---------------------------------------------------------------------------------------------
% Results of the experiments (folder DLPOS_revise_13_10_26): run  python run_all.py  (or every experiment) and
% compile this file: the tables and figures are read directly from the experiment folders (see below).
% A table or figure that does not exist yet is shown as a red "pending" box.
% Two kinds of marks remain, both to be removed before submission:
%   \res{...}          a value known only to the authors
%   \analysisbox{...}  the analysis of the Wilcoxon tests and the conclusion of an experiment, to be
%                      written by the authors once the tables of the experiment are available
% ---------------------------------------------------------------------------------------------
\newcommand{\res}[1]{\textcolor{red}{#1}}
\newcommand{\analysisbox}[1]{\par\medskip\noindent\fcolorbox{blue}{white}{\parbox{\dimexpr\linewidth-2\fboxsep-2\fboxrule\relax}{\small\textcolor{blue}{\textbf{Analysis of the Wilcoxon tests and conclusion (to be written by the authors from the tables above).} #1}}}\par\medskip}
% ---------------------------------------------------------------------------------------------
% Results of the experiments (DLPOS_revise_13_10_26).  Every table (\resulttable{name} -> name.tex) and every
% figure (\resultfig{options}{file}) is read DIRECTLY from the folder of the experiment that produces it,
% so that after running the experiments one compilation updates the whole paper.  The folders are
% searched in the order below; paper/results/ (filled by collect_results.py, e.g. to make a self-contained
% source for the journal) is used last.  Missing files appear as red "Pending" boxes.
% ---------------------------------------------------------------------------------------------
\newcommand{\resultdirs}{../results/}   % figures and tables of the paper (folder paper/results)
\makeatletter
\newcommand{\findresult}[1]{% sets \resultpath to the first folder of \resultdirs that contains #1
  \global\let\resultpath\relax
  \@for\@rd:=\resultdirs\do{%
    \ifx\resultpath\relax\IfFileExists{\@rd#1}{\xdef\resultpath{\@rd#1}}{}\fi}}
\newcommand{\pendingbox}[1]{\fbox{\parbox{0.9\linewidth}{\centering\textcolor{red}{Pending: \detokenize{#1} (not found in the experiment folders)}}}}
\newcommand{\resulttable}[1]{\findresult{#1.tex}%
  \ifx\resultpath\relax\expandafter\@secondoftwo\else\expandafter\@firstoftwo\fi
  {{\let\resultorigtable\table\def\table[##1]{\resultorigtable[!htbp]}\input{\resultpath}}}% generated tables use [H]: placed as floats
  {\begin{table}[!htbp]\centering\pendingbox{#1.tex}\end{table}}}
\newcommand{\resultfig}[2]{\findresult{#2}%
  \ifx\resultpath\relax\expandafter\@secondoftwo\else\expandafter\@firstoftwo\fi
  {\includegraphics[#1]{\resultpath}}{\pendingbox{#2}}}
\makeatother

\begin{document}

\begin{center}
{\Large\bfseries Supplementary material}\\[6pt]
{\large Dynamic Lifecycle Particle Swarm Optimization: A New Swarm-based Metaheuristic Approach}\\[6pt]
Nguyen Ngoc Dang Duy, Kien Trung Nguyen, Nguyen Thanh Hung
\end{center}

\noindent Section numbers (e.g. Section~3.1) refer to the sections of the paper. The tables and figures of this material are numbered A1, A2, \dots

\appendix
% figures and tables of the appendix are numbered A1, A2, ... so that references from the main text show where they are
\setcounter{figure}{0}\setcounter{table}{0}\renewcommand{\thefigure}{A\arabic{figure}}\renewcommand{\thetable}{A\arabic{table}}
\FloatBarrier
\section{Detailed results}\label{app:detailed}
This supplementary material gives the per-function results summarised in Section~3 of the paper, together with the supplementary tables and figures referred to in that section. All values are computed over independent runs under the protocol of Section~3. Every run (best value, best solution, convergence curve, run time, seed and platform) is recorded and available from the authors, so that all statistics can be recomputed.

\FloatBarrier
\subsection{Velocity functions $v_1$ and $v_2$ (Section~3.1)}\label{app:v1v2}
This part supports Section~3.1. Table~\ref{tab:v1_v2_tails} and Figure~\ref{fig:v1_v2_tails} describe the tails of $v_1$ and $v_2$ in the pack (part~(B)). Figure~\ref{fig:v1_v2_cost} shows the time per value of the two functions and Table~\ref{tab:v1_v2_share} the share of the run time of DLPSO spent in the velocity function (part~(A)). Figure~\ref{fig:v1_v2_alpha} shows the mean Friedman ranks for every steepness, and Tables~\ref{tab:v1_v2_D10}--\ref{tab:v1_v2_D10:p3} ($D=10$), Tables~\ref{tab:v1_v2_D30}--\ref{tab:v1_v2_D30:p3} ($D=30$) and Figure~\ref{sosanhham} give the per-function results for the default steepness $\alpha=\frac1{10}$ (part~(C)).


\resulttable{table_v1_v2_tails}


\begin{figure}[!htbp]
    \centering
    \resultfig{width=\textwidth}{v1_v2_tails.pdf}
    \caption{(a) $v_1$ and $v_2$ as functions of $z=(\varphi-\frac12)/\alpha$: the quadratic exponent of $v_1$ gives a parabola, the linear exponent of $v_2$ two straight lines. The vertical lines mark birth and end of life, $|z|=1/(2\alpha)$, for the values of $\alpha$ of Table~\ref{tab:v1_v2_tails}. (b) Mean velocity factor $\bar v(t)$ of the pack of DLPSO (default development types, dashed: $v_1$, solid: $v_2$)}
    \label{fig:v1_v2_tails}
\end{figure}


\begin{figure}[!htbp]
    \centering
    \resultfig{width=0.6\textwidth}{v1_v2_cost.pdf}
    \caption{Time per value of the velocity functions $v_1$ and $v_2$ as a function of the number of values evaluated together}
    \label{fig:v1_v2_cost}
\end{figure}

\begin{figure}[!htbp]
    \centering
    \resultfig{width=0.9\textwidth}{v1_v2_alpha.pdf}
    \caption{Mean Friedman ranks of DLPSO with $v_1$ and with $v_2$ as functions of the steepness of the profile (CEC2017, $D=10$ and $D=30$, lower is better)}
    \label{fig:v1_v2_alpha}
\end{figure}


\begin{figure}[!htbp]
   \centering
   \subfloat[Best value $f(x_{best})$ (30 runs) \label{of}]{%
    \resultfig{width=\textwidth,height=0.36\textheight,keepaspectratio}{Box_Value_v1v2_D10.pdf}}
    \\
   \subfloat[CPU time (30 runs) \label{ct}]{%
    \resultfig{width=\textwidth,height=0.36\textheight,keepaspectratio}{Box_Time_v1v2_D10.pdf}}
\caption{Distribution over the runs of the best value $f(x_{best})$ and of the CPU time (seconds) of DLPSO with the velocity functions $v_1$ and $v_2$ on CEC2017 ($D=10$, $\alpha=\frac1{10}$, box: first quartile $Q_1$ to third quartile $Q_3$, line: median, marker: mean, whiskers: most extreme runs within $1.5\,\mathrm{IQR}$ of the box, points: other runs)}
    \label{sosanhham}
\end{figure}


\resulttable{table_v1_v2_D10}

\resulttable{table_v1_v2_D30}

\resulttable{table_v1_v2_share}

\FloatBarrier
\subsection{Sensitivity analysis (Section~3.2)}\label{app:sens}
This part supports Section~3.2. Tables~\ref{tab:sensitivity}--\ref{tab:sensitivity:p4} give, for every value of every parameter, the mean Friedman rank, the counts $+/=/-$ of the default and the Holm-adjusted $p$-value of the signed-rank test. Tables~\ref{tab:sens_convergence}--\ref{tab:sens_convergence:p3} and Figures~\ref{fig:sens_convergence}--\ref{fig:sens_convergence_4} describe the convergence speed for every value of every parameter.


\resulttable{table_sensitivity}

\resulttable{table_sens_convergence}

\begin{figure}[!htbp]
    \centering
    \resultfig{width=\textwidth}{sens_convergence_1.pdf}
    \caption{Normalised convergence progress of DLPSO for every value of the parameters $\eta$, $s$, $\Delta_b$, $(P_1,P_2)$ (rows) on the classical functions and on CEC2017 with $D=10$ and $D=30$ (columns), averaged over the functions of the suite and over 30 runs (red: default, 1: best final value among the values of the parameter) (part 1/4)}
    \label{fig:sens_convergence}
\end{figure}

\begin{figure}[!htbp]
    \centering
    \resultfig{width=\textwidth}{sens_convergence_2.pdf}
    \caption{Normalised convergence progress of DLPSO for every value of the parameters $\alpha$, $A_{\max}$, $\rho$, $\delta$ (rows) on the classical functions and on CEC2017 with $D=10$ and $D=30$ (columns), averaged over the functions of the suite and over 30 runs (red: default, 1: best final value among the values of the parameter) (part 2/4)}
    \label{fig:sens_convergence_2}
\end{figure}

\begin{figure}[!htbp]
    \centering
    \resultfig{width=\textwidth}{sens_convergence_3.pdf}
    \caption{Normalised convergence progress of DLPSO for every value of the parameters noise width $c_r$, $v_{\max}$, $\kappa_0$, $H$ (rows) on the classical functions and on CEC2017 with $D=10$ and $D=30$ (columns), averaged over the functions of the suite and over 30 runs (red: default, 1: best final value among the values of the parameter) (part 3/4)}
    \label{fig:sens_convergence_3}
\end{figure}

\begin{figure}[!htbp]
    \centering
    \resultfig{width=\textwidth}{sens_convergence_4.pdf}
    \caption{Normalised convergence progress of DLPSO for every value of the parameters $\sigma_\kappa$, $\mu_i^t$, $\varepsilon$, $N$ (rows) on the classical functions and on CEC2017 with $D=10$ and $D=30$ (columns), averaged over the functions of the suite and over 30 runs (red: default, 1: best final value among the values of the parameter) (part 4/4)}
    \label{fig:sens_convergence_4}
\end{figure}

\FloatBarrier
\subsection{Convergence behaviour of DLPSO (Section~3.3)}\label{app:conv}
This part supports Section~3.3. Figures~\ref{fig:conv_classical} and \ref{fig:conv_cec10} show the convergence of DLPSO on the classical functions and on CEC2017 with $D=10$, and Tables~\ref{tab:wilcoxon_convergence_functions}--\ref{tab:wilcoxon_convergence_functions:p3} give the Wilcoxon tests of the improvement of DLPSO on every function.


\begin{figure}[!htbp]
    \centering
    \resultfig{width=\textwidth,height=0.42\textheight,keepaspectratio}{Conv_DLPSO_classical.pdf}
    \caption{Convergence of DLPSO on the classical functions F1--F18 ($D=10$, $N=50$, 30 runs): mean optimality gap $f(x_{best})-f^*$ as a function of the fraction of MaxFE used, with the band between the first and the third quartile of the runs}
    \label{fig:conv_classical}
\end{figure}


\resulttable{table_wilcoxon_convergence_functions}

\begin{figure}[!htbp]
    \centering
    \resultfig{width=\textwidth,height=0.42\textheight,keepaspectratio}{Conv_DLPSO_CEC_D10.pdf}
    \caption{Convergence of DLPSO on CEC2017 ($D=10$, $N=100$, 30 runs): mean distance to the optimum $f(x_{best})-F^*$ with the band between the first and the third quartile}
    \label{fig:conv_cec10}
\end{figure}

\FloatBarrier
\subsection{Classical benchmark functions (Section~3.4)}\label{app:low}
This part supports Section~3.4. Tables~\ref{3tt}--\ref{3tt:p8} give the per-function results of the five algorithms for every population size. Figure~\ref{benchmarkop} shows the distributions of the optimality gaps with 10 and 20 particles, Figures~\ref{benchmarktime} and \ref{benchmarktime2} the distributions of the CPU times, and Figures~\ref{fig:conv_low_10}--\ref{fig:conv_low_50} the convergence curves for 10, 20, 30 and 50 particles.


\begin{figure}[!htbp]
    \centering
    \subfloat[10 particles \label{10con}]{%
       \resultfig{width=\textwidth}{Box_Gap_Pop_10.pdf}}
    \\[0.2cm]
    \subfloat[20 particles \label{20con}]{%
       \resultfig{width=\textwidth}{Box_Gap_Pop_20.pdf}}
    \caption{Distribution over the runs of the optimality gap $|f(x_{best})-f^*|$ on the classical functions with 10 and 20 particles (same budget, 30 runs, box: first quartile $Q_1$ to third quartile $Q_3$, line: median, marker: mean, whiskers: most extreme runs within $1.5\,\mathrm{IQR}$ of the box, points: other runs)}
    \label{benchmarkop}
\end{figure}

\begin{figure}[!htbp]
    \centering
    \subfloat[10 particles \label{10con-time}]{%
       \resultfig{width=\textwidth}{Box_Time_Pop_10.pdf}}
    \\[0.2cm]
    \subfloat[20 particles \label{20con-time}]{%
       \resultfig{width=\textwidth}{Box_Time_Pop_20.pdf}}
    \caption{Distribution over the runs of the CPU time (seconds) on the classical functions with 10 and 20 particles (same budget, 30 runs, box: first quartile $Q_1$ to third quartile $Q_3$, line: median, marker: mean, whiskers: most extreme runs within $1.5\,\mathrm{IQR}$ of the box, points: other runs)}
    \label{benchmarktime}
\end{figure}

\begin{figure}[!htbp]
    \centering
    \subfloat[30 particles \label{30con-time}]{%
       \resultfig{width=\textwidth}{Box_Time_Pop_30.pdf}}
    \\[0.2cm]
    \subfloat[50 particles \label{50con-time}]{%
       \resultfig{width=\textwidth}{Box_Time_Pop_50.pdf}}
    \caption{Distribution over the runs of the CPU time (seconds) on the classical functions with 30 and 50 particles (same budget, 30 runs, box: first quartile $Q_1$ to third quartile $Q_3$, line: median, marker: mean, whiskers: most extreme runs within $1.5\,\mathrm{IQR}$ of the box, points: other runs)}
    \label{benchmarktime2}
\end{figure}

\begin{figure}[!htbp]
    \centering
    \resultfig{width=\textwidth,height=0.42\textheight,keepaspectratio}{Convergence_Pop_10.pdf}
    \caption{Convergence curves of DLPSO, PSO, BA, DE and CMA-ES on all the classical functions F1--F18 with 10 particles (mean over 30 runs of the best-so-far optimality gap $|f(x_{best})-f^*|$ as a function of the fraction of MaxFE used)}
    \label{fig:conv_low_10}
\end{figure}

\begin{figure}[!htbp]
    \centering
    \resultfig{width=\textwidth,height=0.42\textheight,keepaspectratio}{Convergence_Pop_20.pdf}
    \caption{Convergence curves of DLPSO, PSO, BA, DE and CMA-ES on all the classical functions F1--F18 with 20 particles (mean over 30 runs of the best-so-far optimality gap $|f(x_{best})-f^*|$ as a function of the fraction of MaxFE used)}
    \label{fig:conv_low_20}
\end{figure}

\begin{figure}[!htbp]
    \centering
    \resultfig{width=\textwidth,height=0.42\textheight,keepaspectratio}{Convergence_Pop_30.pdf}
    \caption{Convergence curves of DLPSO, PSO, BA, DE and CMA-ES on all the classical functions F1--F18 with 30 particles (mean over 30 runs of the best-so-far optimality gap $|f(x_{best})-f^*|$ as a function of the fraction of MaxFE used)}
    \label{fig:conv_low_30}
\end{figure}

\begin{figure}[!htbp]
    \centering
    \resultfig{width=\textwidth,height=0.42\textheight,keepaspectratio}{Convergence_Pop_50.pdf}
    \caption{Convergence curves of DLPSO, PSO, BA, DE and CMA-ES on all the classical functions F1--F18 with 50 particles (mean over 30 runs of the best-so-far optimality gap $|f(x_{best})-f^*|$ as a function of the fraction of MaxFE used)}
    \label{fig:conv_low_50}
\end{figure}


\resulttable{table_low_dim}

\FloatBarrier
\subsection{CEC2017, $D=10$ and $D=30$ (Section~3.5)}\label{app:cec}
This part supports Section~3.5. Tables~\ref{cec2017_D10}--\ref{cec2017_D10:p4} ($D=10$) and \ref{cec2017_D30}--\ref{cec2017_D30:p4} ($D=30$) give the per-function results. Figures~\ref{cec2017_op} and \ref{cec2017_conv10} show the distributions of the best values and the convergence curves for $D=10$, and Figure~\ref{cec2017_time} the distributions of the CPU times for $D=30$.


\begin{figure}[!htbp]
    \centering
        \resultfig{width=\textwidth}{Box_Value_D10.pdf}
    \caption{Distribution over 30 runs of the best value $f(x_{best})$ of DLPSO, PSO, BA, DE and CMA-ES on the CEC2017 functions ($D=10$, box: first quartile $Q_1$ to third quartile $Q_3$, line: median, marker: mean, whiskers: most extreme runs within $1.5\,\mathrm{IQR}$ of the box, points: other runs)}
    \label{cec2017_op}
\end{figure}


\begin{figure}[!htbp]
    \centering
    \resultfig{width=\textwidth}{Convergence_D10.pdf}
    \caption{Convergence curves of DLPSO, PSO, BA, DE and CMA-ES on all the CEC2017 functions ($D=10$, mean over 30 runs of the best-so-far value $f(x_{best})$ as a function of the number of function evaluations)}
    \label{cec2017_conv10}
\end{figure}

\begin{figure}[!htbp]
    \centering
        \resultfig{width=\textwidth}{Box_Time_D30.pdf}
        \caption{Distribution of the CPU time (seconds, 30 runs) of one run of DLPSO, PSO, BA, DE and CMA-ES on the CEC2017 functions ($D=30$, MaxFE $=3\cdot10^5$, box: first quartile $Q_1$ to third quartile $Q_3$, line: median, marker: mean, whiskers: most extreme runs within $1.5\,\mathrm{IQR}$ of the box, points: other runs)}
    \label{cec2017_time}
\end{figure}


\resulttable{table_cec2017_D10}

\resulttable{table_cec2017_D30}

\FloatBarrier
\subsection{CEC2017, $D=50$ and $D=100$ (Section~3.6)}\label{app:high}
This part supports Section~3.6. Tables~\ref{cec2017_D50}--\ref{cec2017_D50:p4} ($D=50$) and \ref{cec2017_D100}--\ref{cec2017_D100:p4} ($D=100$) give the per-function results. Figures~\ref{fig:convergence_50D} and \ref{fig:convergence_100D} show the convergence curves, and Figure~\ref{fig:time_100D} the distributions of the CPU times for $D=100$.


\begin{figure}[!htbp]
    \centering
    \resultfig{width=\textwidth}{Convergence_D50.pdf}
    \caption{Convergence curves of DLPSO, PSO, BA, DE and CMA-ES on all the CEC2017 functions ($D=50$, mean over 30 runs of the best-so-far value $f(x_{best})$)}
    \label{fig:convergence_50D}
\end{figure}

\begin{figure}[!htbp]
    \centering
    \resultfig{width=\textwidth}{Convergence_D100.pdf}
    \caption{Convergence curves of DLPSO, PSO, BA, DE and CMA-ES on all the CEC2017 functions ($D=100$, mean over 30 runs of the best-so-far value $f(x_{best})$)}
    \label{fig:convergence_100D}
\end{figure}

\begin{figure}[!htbp]
    \centering
    \resultfig{width=\textwidth}{Box_Time_D100.pdf}
    \caption{Distribution of the CPU time (seconds, 30 runs) of one run on the CEC2017 functions ($D=100$, MaxFE $=10^6$, box: first quartile $Q_1$ to third quartile $Q_3$, line: median, marker: mean, whiskers: most extreme runs within $1.5\,\mathrm{IQR}$ of the box, points: other runs)}
    \label{fig:time_100D}
\end{figure}


\resulttable{table_cec2017_D50}

\resulttable{table_cec2017_D100}

\FloatBarrier
\subsection{CEC2013 LSGO, $D=1000$ (Section~3.7)}\label{app:lsgo}
This part supports Section~3.7. Tables~\ref{lsgo2013_D1000}--\ref{lsgo2013_D1000:p2} give the per-function results of the four algorithms, and Figure~\ref{fig:time_lsgo} shows the distributions of their CPU times.


\begin{figure}[!htbp]
    \centering
    \resultfig{width=\textwidth}{Box_Time_D1000.pdf}
    \caption{Distribution of the CPU time (seconds, 30 runs) of the four algorithms on the CEC2013 LSGO functions ($D=1000$, MaxFE $=1.2\cdot10^5$, box: first quartile $Q_1$ to third quartile $Q_3$, line: median, marker: mean, whiskers: most extreme runs within $1.5\,\mathrm{IQR}$ of the box, points: other runs)}
    \label{fig:time_lsgo}
\end{figure}


\resulttable{table_lsgo2013_D1000}


\FloatBarrier
\end{document}
